\section{Semi-Structured Interview Guide for Expert Interviews with Facilitators}
\label{annex:interview-guide-facilitators}

This annex documents the original-language instrument (German) used in the expert interviews
(Chapter~\ref{ch:empirical-analysis-of-facilitating-strategies}). The guide operationalizes
RQ1 (facilitating strategies and challenges in goal-oriented remote meetings with convergent
objectives) and prepares the derivation of design requirements for RQ2 (a proactive,
context-sensitive AI moderation system). It does not address RQ3 (user acceptance among
knowledge workers), which is examined in a separate study.

The instrument follows the logic of a semi-structured expert interview: a small set of open
narrative prompts (\textit{Hauptfragen}) secures thematic focus, while optional probes
(\textit{Nachfragen}) are used only if a theoretically relevant aspect is not addressed
spontaneously \cite[p.~877]{baurHandbuchMethodenEmpirischen2022}. The sequence may follow
the conversation; not every probe needs to be asked. The intended duration is approximately
30 minutes.

\subsection*{Expected thematic clusters and coverage}

The blocks are aligned with the challenges identified in the literature
(Chapter~2) and with the system-design dimensions of RQ2 (detection, timing, addressee,
and form of intervention). The following clusters are expected in the material; they serve as
a sensitizing frame for later thematic coding, not as a closed coding scheme.

\begin{enumerate}[label=\textbf{C\arabic*}]
  \item \textbf{Goal clarity and structure.}
    Unclear or poorly communicated objectives; agenda without operational sub-goals;
    loss of a shared sense of ``on track''.
    Covered by Block~2 and Block~3.
  \item \textbf{Thematic deviation and time-to-decision.}
    Off-topic sequences, circular discussion, delayed or avoided decisions in convergent
    processes (decision-making, prioritization, problem-solving).
    Covered by Block~2 and Block~4.
  \item \textbf{Participation, hierarchy, and group dynamics.}
    Unequal speaking time, dominance, reluctance to dissent, face-saving.
    Covered by Block~2, Block~4 (public vs.\ private), and Block~6.
  \item \textbf{Dynamic goals and situational uncertainty.}
    Goals that shift during the meeting; false assumptions; need to renegotiate the mandate.
    Covered by Block~2 and Block~3.
  \item \textbf{Preparation versus in-meeting facilitation.}
    What can be secured beforehand, and what can only be steered in situ.
    Covered by Block~3.
  \item \textbf{Recognizing critical moments (cues for detection).}
    Verbal, temporal, social, and remote-specific signals from which facilitators infer
    that goal orientation or decision quality is at risk. Directly informs the detection
    logic of RQ2 (\textit{What? When? Who?}).
    Covered by Block~4.
  \item \textbf{Intervention repertoire.}
    Methods, structures, and habits (e.g.\ timeboxing, roles, visualization);
    choice of addressee and channel (group vs.\ private); tone and timing.
    Covered by Block~3 and Block~4.
  \item \textbf{Tool support and gaps.}
    How current platforms help or hinder goal orientation; missing support.
    Covered by Block~5.
  \item \textbf{Role, agency, and limits of an AI co-facilitator.}
    Proactive versus on-demand support; essential functions; undesired automation;
    risks of unsolicited intervention (trust, control, false positives, privacy).
    Covered by Block~6.
\end{enumerate}

Removed from the earlier draft were biographical and attitudinal items (general opinion of
  online collaboration, career entry, liking the role, overall rating of online vs.\ in-person
work). They do not contribute to RQ1 or RQ2 and consume time needed for narrative accounts
of practice.

\subsection*{Hinweise zur Durchführung}

\begin{itemize}[leftmargin=*]
  \item Anrede durchgängig \textit{Sie}; Fachbegriffe der interviewten Person aufnehmen
    (Moderation, Facilitation, Prozessbegleitung).
  \item Hauptfragen als Erzählaufforderungen stellen, nicht als Ja/Nein- oder
    Bewertungsfragen.
  \item Immanente Nachfragen (am Gesagten anknüpfen) haben Vorrang vor exmanenten Sondierungen.
  \item Die optionalen Nachfragen sind Sensibilisierungsanker aus der Literatur; sie sollen
    nicht als Checkliste abgefragt werden.
  \item Mit \textit{Entscheidung} ist in diesem Leitfaden ein verbindliches Gruppenergebnis
    gemeint, nicht beliebige Meinungsäußerung: die Gruppe legt fest, wie es weitergeht,
    was priorisiert wird oder worauf sie sich als Team committet (Ziele, nächste Schritte,
    gemeinsame Arbeitsweisen). Diese Arbeitsdefinition bei Bedarf im Gespräch wiederholen.
  \item Nach dem Gespräch: unmittelbares Interviewprotokoll (erste Eindrücke, besondere
    Sequenzen, Kontext) anfertigen, solange der Eindruck noch frisch ist.
\end{itemize}

\subsection*{Einstiegstext (sinngemäß)}

Vielen Dank, dass Sie sich Zeit nehmen. Ich untersuche, wie zielorientierte Remote-Meetings
so moderiert werden, dass die Gruppe beim vereinbarten Ziel bleibt und zu einem
verbindlichen Ergebnis kommt, also zu einer Entscheidung im engeren Sinn.

Gemeint sind damit Entscheidungssituationen, in denen die Gruppe zu einer Festlegung gelangen
muss \emph{wie es weitergeht}. Die Art der Enscheidung kann dabei sehr vielfälltig sein. Typische Fälle sind
zum Beispiel:
\begin{itemize}
  \item Priorisierung, etwa welches Vorhaben oder Feature als Nächstes umgesetzt wird.
  \item Commitment des Teams, etwa welche Ziele es sich setzt.
  \item verbindliche nächste Schritte oder die Festlegung gemeinsamer Arbeitsweisen.
\end{itemize}
Entscheidend ist, dass aus dem Meeting eine gemeinsame Festlegung folgt, an der sich die
weitere Arbeit ausrichten kann.

Aus den Gesprächen mit erfahrenen Moderatorinnen und Moderatoren leite ich ab, welche
Anforderungen ein KI-gestütztes Unterstützungssystem erfüllen müsste, um Teilnehmende dabei
zu unterstützen, genau bei solchen Aufträgen am Ziel zu bleiben und zügiger zu einer
solchen Festlegung zu kommen.

Das Gespräch ist als offenes Experteninterview angelegt. Es gibt keine richtigen oder
falschen Antworten; mich interessiert Ihre berufliche Praxis und Ihre Einschätzung. Das
Interview dauert etwa 30 Minuten, wird über Zoom geführt und --- mit Ihrem Einverständnis ---
aufgezeichnet, transkribiert und in anonymisierter Form für die Bachelorarbeit ausgewertet.
Die Aufnahme wird nicht an Dritte weitergegeben. Sind Sie mit Aufzeichnung und dieser
Verarbeitung einverstanden?

\medskip
\textit{Einverständnis dokumentieren, dann aufzeichnen.}

% ---------------------------------------------------------------------------
\subsection*{Block 1: Beruflicher Kontext}
\textit{Funktion: Sampling-Validierung und gemeinsame Bezugsebene. Ca.\ 3 Minuten.}

\begin{enumerate}[label=1.\arabic*, leftmargin=*]
  \item Bitte beschreiben Sie kurz Ihre Rolle und Ihren Erfahrungshintergrund in der
    Moderation --- insbesondere in Remote-Meetings, in denen die Gruppe zu einer
    verbindlichen Festlegung kommen soll (nächster Schritt, Priorisierung, Ziele oder
    gemeinsame Arbeitsweisen).
    \begin{itemize}[label=\textendash]
      \item \textit{Optional:} Seit wann sind Sie in diesem Feld tätig, und in welchen
        organisatorischen Kontexten (z.\,B.\ Unternehmen, öffentliche Verwaltung,
        agiles Coaching, Großgruppen)?
      \item \textit{Optional:} Wer sind typischerweise die Teilnehmenden (Hierarchie,
        Fachlichkeit, Gruppengröße, wiederkehrendes Team vs.\ ad hoc)?
    \end{itemize}
\end{enumerate}

% ---------------------------------------------------------------------------
\subsection*{Block 2: Herausforderungen in zielorientierten Remote-Meetings (RQ1)}
\textit{Funktion: Empirische Verdichtung der in Kapitel~2 identifizierten Herausforderungen;
  Fokus auf konvergente Prozesse, nicht auf Online-Zusammenarbeit im Allgemeinen.
Ca.\ 6 Minuten.}

\begin{enumerate}[label=2.\arabic*, leftmargin=*]
  \item Wenn Sie an typische Remote-Meetings denken, in denen eine Entscheidung getroffen
    werden soll: Was läuft dort aus Ihrer Erfahrung häufig schief --- und woran merken
    Sie das im Verlauf des Gesprächs?
    \begin{itemize}[label=\textendash]
      \item \textit{Falls nicht von selbst thematisiert:} Zielklarheit und Struktur;
        Abschweifen vom Thema; ungleiche Beteiligung oder hierarchische Dynamiken;
        Zeitdruck; Ziele, die sich während des Meetings verschieben.
    \end{itemize}
  \item Welche dieser Schwierigkeiten sind für Sie spezifisch für das Remote-Setting ---
    und welche würden in vergleichbarer Form auch in Präsenz auftreten?
\end{enumerate}

% ---------------------------------------------------------------------------
\subsection*{Block 3: Strategien der Ziel- und Entscheidungssicherung (RQ1)}
\textit{Funktion: Rekonstruktion bewährter Praktiken (Methoden, Strukturen, Gewohnheiten)
und der Sichtbarmachung des Ziels für die Gruppe. Ca.\ 6 Minuten.}

\begin{enumerate}[label=3.\arabic*, leftmargin=*]
  \item Wie gehen Sie vor, damit ein solches Meeting am Ziel bleibt und die Gruppe zu einer
    Entscheidung kommt? Beschreiben Sie gerne ein konkretes Vorgehen anhand einer
    typischen Situation.
    \begin{itemize}[label=\textendash]
      \item \textit{Optional:} Welche Methoden, Strukturen oder Gewohnheiten haben sich
        für Sie bewährt (z.\,B.\ Timeboxing, Rollen, Visualisierung, Entscheidungsregeln)?
      \item \textit{Optional:} Was davon legen Sie in der Vorbereitung fest, und was
        steuern Sie erst in der Situation?
    \end{itemize}
  \item Wie stellen Sie sicher, dass den Teilnehmenden während des Meetings klar bleibt,
    was in diesem Moment ``on track'' bedeutet --- also was das aktuelle Ziel bzw.\ der
    aktuelle Entscheidungsauftrag ist?
    \begin{itemize}[label=\textendash]
      \item \textit{Optional:} Was geschieht, wenn das Ziel selbst unklar ist oder sich
        im Meeting als unpassend erweist?
    \end{itemize}
\end{enumerate}

% ---------------------------------------------------------------------------
\subsection*{Block 4: Kritische Momente, Erkennung und Intervention (RQ1, Anschluss an RQ2)}
\textit{Funktion: Explikation impliziten Expertenwissens über Zeitpunkt, Anlass und Form
  von Interventionen. Die Cues sind zentral für die spätere Spezifikation der
Erkennungslogik. Ca.\ 7 Minuten.}

\begin{enumerate}[label=4.\arabic*, leftmargin=*]
  \item Was sind in zielorientierten Meetings mit Entscheidungsauftrag typische kritische
    Momente? Was kann dort konkret schiefgehen --- für das Ziel, für die Entscheidung
    und für die Gruppe?
  \item Woran erkennen Sie diese Momente? Können Sie das möglichst an konkreten
    Beobachtungen festmachen --- etwa daran, wie gesprochen wird, wer spricht oder
    schweigt, wie die Zeit vergeht, oder woran Sie im Remote-Setting merken, dass die
    Gruppe den Faden verliert?
    \begin{itemize}[label=\textendash]
      \item \textit{Optional, nicht als Messzwang formulieren:} Gibt es Anzeichen, die
        sich für Sie vergleichsweise zuverlässig wiederholen, sodass Sie fast schon
        sagen könnten: Jetzt muss ich eingreifen?
    \end{itemize}
  \item Wann sprechen Sie jemanden direkt in der Gruppe an, wann eher auf einem privaten
    Kanal --- und nach welchen Kriterien treffen Sie diese Wahl?
    \begin{itemize}[label=\textendash]
      \item \textit{Optional:} Was riskieren Sie, wenn Sie zu früh, zu spät, zu öffentlich
        oder zu privat eingreifen?
    \end{itemize}
\end{enumerate}

% ---------------------------------------------------------------------------
\subsection*{Block 5: Werkzeuge und Unterstützungsbedarf (Brücke RQ1 $\rightarrow$ RQ2)}
\textit{Funktion: Bestehende sozio-technische Praxis und Lücken, ohne bereits eine
KI-Lösung nahezulegen. Ca.\ 4 Minuten.}

\begin{enumerate}[label=5.\arabic*, leftmargin=*]
  \item Welche Plattformen und Werkzeuge nutzen Sie in der Moderation --- und inwiefern
    helfen oder behindern diese Sie dabei, Zielorientierung und Entscheidungsfortschritt
    zu sichern?
    \begin{itemize}[label=\textendash]
      \item \textit{Optional:} Welche Unterstützung hätten Sie sich in konkreten
        Situationen gewünscht, die heutige Werkzeuge nicht leisten?
    \end{itemize}
\end{enumerate}

% ---------------------------------------------------------------------------
\subsection*{Block 6: Anforderungen an ein KI-gestütztes Moderationssystem (RQ2)}
\textit{Funktion: Offene Anforderungserhebung aus Expertensicht. Das System wird als
  Unterstützung der Teilnehmenden bei Zielbindung und Entscheidungsgeschwindigkeit
eingeführt, nicht als Ersatz der menschlichen Moderation. Ca.\ 7 Minuten.}

\begin{enumerate}[label=6.\arabic*, leftmargin=*]
  \item Stellen Sie sich ein KI-System vor, das in Remote-Meetings die Teilnehmenden
    dabei unterstützt, am vereinbarten Ziel zu bleiben und zügiger zu einer Entscheidung
    zu kommen. Wie sollte ein solches System aus Ihrer Sicht agieren --- und was sollte
    es ausdrücklich nicht tun?
    \begin{itemize}[label=\textendash]
      \item \textit{Optional:} Eher proaktiv von sich aus, oder nur auf Anfrage?
      \item \textit{Optional:} Welche Rolle wäre angemessen (z.\,B.\ stiller Spiegel,
        Zeit- und Zielwächter, Co-Moderatorin, Vorschlagsgeberin)?
    \end{itemize}
  \item Welche Funktionen wären für Sie essenziell, damit das System tatsächlich nützt ---
    bezogen auf das Erkennen von Abweichungen, auf den Zeitpunkt und die Form der
    Intervention sowie darauf, wen das System anspricht (Gruppe oder einzelne Person)?
    \begin{itemize}[label=\textendash]
      \item \textit{Optional:} Wenn das System zwischen öffentlicher Intervention und privater
        Rückmeldung wählen könnte: Nach welchen Kriterien sollte es das tun?
    \end{itemize}
  \item Welche Risiken, Schwächen oder unbeabsichtigten Wirkungen sehen Sie, wenn ein
    KI-System von sich aus in Meetings einschreitet?
    \begin{itemize}[label=\textendash]
      \item \textit{Falls nicht von selbst thematisiert:} Vertrauen und Kontrolle der
        Gruppe; Fehlalarme; Gesichtswahrung; hierarchische Verzerrung; Datenschutz und
        Transparenz.
    \end{itemize}
\end{enumerate}

% ---------------------------------------------------------------------------
\subsection*{Abschluss}

\begin{enumerate}[label=7.\arabic*, leftmargin=*]
  \item Gibt es Aspekte der Zielorientierung, der Entscheidungsfindung oder der
    KI-Unterstützung, die wir nicht angesprochen haben und die Ihnen wichtig erscheinen?
\end{enumerate}

\medskip
Dank, Hinweis auf Anonymisierung und --- falls zugesagt --- auf die Möglichkeit, Ergebnisse
zurückzuspiegeln. Aufnahme beenden.
